\chapter{Geodesics and the Carter Constant}
\label{ch:geodesics}

\section{Constants of motion and the Carter separation}
\label{sec:carter}

Geodesics in Kerr admit four constants of motion: the energy $E$, the axial
angular momentum $L_z$, the Carter constant $Q$ \citep{Carter1968}, and the
norm $\mu^2$ of the four-momentum. With
$P(r) = E(r^2+a^2)-aL_z$ the Carter-separated equations for a photon
($\mu=0$) read
\begin{align}
\Sigma\,\frac{\mathrm{d}r}{\mathrm{d}\lambda} &= \pm\sqrt{R(r)},\nonumber\\
\Sigma\,\frac{\mathrm{d}\theta}{\mathrm{d}\lambda} &=
  \pm\sqrt{\Theta(\theta)},\nonumber\\
\Sigma\,\frac{\mathrm{d}\phi}{\mathrm{d}\lambda} &=
  -\Big(aE - \frac{L_z}{\sin^2\theta}\Big) + \frac{a}{\Delta}P(r),
  \nonumber\\
\Sigma\,\frac{\mathrm{d}t}{\mathrm{d}\lambda} &=
  -a(aE\sin^2\theta - L_z) + \frac{(r^2+a^2)}{\Delta}P(r),
\label{eq:carter}
\end{align}
with the potentials
\begin{align}
R(r) &= P(r)^2 - \Delta\left[(L_z-aE)^2 + Q\right],
\label{eq:rpot}\\
\Theta(\theta) &= Q - \cos^2\theta\left[a^2(-E^2) +
\frac{L_z^2}{\sin^2\theta}\right].
\label{eq:thetapot}
\end{align}
The potentials and their derivatives are implemented by
\texttt{radial\_potential}, \texttt{polar\_potential} and their
\texttt{\_deriv} companions in \texttt{physics/src/null\_geodesics.cpp}.

\stub{The Hamilton--Jacobi derivation of the separation; $Q$ from the
Killing--Yano tensor; turning points and orbit classification by the roots
of $R(r)$. New figures: sample photon and timelike trajectories;
turning-point structure of $R(r)$.}

\section{Mino time and the second-order formulation}
\label{sec:mino}

Rather than integrating the first-order radial and polar equations with
explicit sign bookkeeping at turning points, we adopt Mino time
$\tau$ \citep{Mino2003}, defined by
$\mathrm{d}\lambda = \Sigma\,\mathrm{d}\tau$, which decouples the $(r)$ and
$(\theta)$ sectors, and evolve the \emph{second-order} equations obtained
by differentiating $\dot r^2 = R(r)$ once:
\begin{equation}
\frac{\mathrm{d}^2 r}{\mathrm{d}\tau^2} = \tfrac{1}{2}R'(r), \qquad
\frac{\mathrm{d}^2 \theta}{\mathrm{d}\tau^2} = \tfrac{1}{2}\Theta'(\theta),
\label{eq:secondorder}
\end{equation}
in which turning points are ordinary smooth extrema and no square root of a
near-zero quantity is ever taken. The first integrals
$\dot r^2 = R(r)$ and $\dot\theta^2 = \Theta(\theta)$ then serve as exact
algebraic constraints on the numerical solution, which we re-impose after
every accepted step (Chapter~\ref{ch:numerics}). The coordinate
singularities at $\theta = 0,\pi$ are regularised with a $10^{-8}$ clamp;
they are properties of the BL chart, not of the spacetime.

\section{Timelike geodesics}
\label{sec:timelike}

For timelike geodesics ($\mu=1$, stellar orbits) the radial potential gains
the mass term, $R(r) = P^2 - \Delta[r^2 + (L_z-aE)^2 + Q]$, the polar
potential becomes $\Theta = Q - \cos^2\theta[a^2(1-E^2)+L_z^2/\sin^2\theta]$,
and the four-velocity normalisation $g_{\mu\nu}u^\mu u^\nu = -1$ provides an
independent per-step health check (\texttt{solve\_ut} in
\texttt{timelike\_geodesics.cpp} solves the normalisation quadratic;
\texttt{carter\_q} recovers $Q$ from the state). We integrate stellar
orbits in the full geodesic equations rather than any post-Newtonian
expansion: at the S2 pericentre ($\sim$1470 Schwarzschild radii, orbital
speed 2.6 per cent of $c$) the leading PN corrections are individually at
the $10^{-4}$ level --- the same order as the measured redshift excess
itself --- so the expansion error budget would have to be argued term by
term, whereas the geodesic integrator's error budget is simply its
conservation diagnostics.

\stub{[OPTIONAL] The exact closed-form solution in elliptic integrals /
Weierstrass $\wp$ (Rauch--Blandford, Gralla--Lupsasca) as an analytic
integrator check; Mino-time vs affine-parameter figure.}
